\chapter*{はじめに}
\addcontentsline{toc}{chapter}{はじめに}
\markboth{はじめに}{}

本書は、宇宙流体シミュレーションや重力多体系シミュレーションに進む学部生が、数値計算法の基礎と検証の考え方を学ぶための準備教材である。
B. Ryden ``\textit{Celestial and Stellar Dynamics}'' (2025) の輪講と並行して用い、そこで学ぶ質量モデル・重力場・恒星軌道を数値計算の題材とする。
Ryden の物理内容の説明を置き換えるのではなく、その式から実際に数値を求め、誤差と収束を検証する方法を学ぶ。

宇宙流体・重力多体系シミュレーションでは、物理量の離散化、重力場の計算、時間積分などを組み合わせる。
本書では、それらを最初から同時に扱わず、単純な問題で一つずつ学ぶ。
本文では基本的な計算法と検証方法を扱い、脚注や補足では、今学んでいる計算が将来のシミュレーションのどの処理に相当するかを示す。

\section*{本書のねらい}

数値計算では、コードがエラーなく最後まで実行され、もっともらしい図が得られても、その答えが正しいとは限らない。
格子幅や時間刻みによる離散化誤差、丸め誤差、計算領域や境界条件の設定による誤差が、気づかないうちに結果に入り込むからである。
本書で最も身につけてほしいのは、個々の公式よりも、得られた数値解が信頼できるかを自分で確かめる習慣である。
そのために、本書では次の3つの確認を繰り返し行う。

\begin{itemize}
  \item \textbf{解析解との比較}：解析解が分かっている問題で、コードと数値計算法が正しく働くことを確かめる。
  \item \textbf{保存量による診断}：エネルギーや角運動量など、物理的に保存されるべき量がどの程度保たれているかを調べる。
  \item \textbf{分解能を変えた収束の確認}：格子幅や時間刻みを変えて計算し、答えが収束するか、誤差が期待どおりの次数で減るかを調べる。
\end{itemize}

実際の研究で解く問題の多くには解析解がない。
だからこそ、まず解析解のある問題で手法を検証し、そのうえで保存量と収束を組み合わせて結果を確かめることが重要になる。
また、誤差が見つかったときには、それが離散化、丸め誤差、有限の計算領域や境界条件、あるいは近似したモデルそのもののどれに由来するのかを区別して考える。

\section*{本書の構成}

本編は第1--3章で完結する。第1・2章では Ryden §6.4 の球対称系を題材に数値積分と有限差分法を学び、第3章では §7.1 の球対称ポテンシャル中の軌道を題材に時間積分法を学ぶ。
本編は球対称系のみを扱うが、付録A・B・Cは、それぞれ第1・2・3章で学んだ方法を軸対称系や非軸対称系へ広げる発展内容であり、関心に応じて選んでよい。また、連立一次方程式の解法を詳しく知りたい場合は、数値線形代数の補足である付録Dを参照してほしい。
各 Notebook は単独で先頭から実行でき、本文との対応表と実行順序を冒頭に示している。

{\footnotesize
\begin{center}
\setlength{\tabcolsep}{3.5pt}
\begin{tabularx}{\linewidth}{>{\raggedright\arraybackslash}p{0.14\linewidth}>{\raggedright\arraybackslash}p{0.16\linewidth}>{\raggedright\arraybackslash}p{0.27\linewidth}X}
\toprule
本書 & Rydenテキストとの対応 & 物理的な題材 & 本書で学ぶ数値計算法 \\
\midrule
第1章\newline Notebook 1 & §6.4\newline 特に §6.4.2--6.4.3 & Plummer・Hernquist・NFW などの球対称密度分布 & 数値積分、誤差、収束次数 \\
第2章\newline Notebook 2 & §6.4 & 球対称密度分布からポテンシャル・重力場を求める問題 & 1次元 Poisson 境界値問題、有限差分、数値微分 \\
第3章\newline Notebook 3 & §7.1 & 球対称ポテンシャル中の恒星軌道 & 常微分方程式の初期値問題、Euler 法、古典的4次 Runge--Kutta 法（RK4）、リープフロッグ法 \\
\midrule
付録A\newline Notebook A & §6.5 & 薄い軸対称円盤、指数円盤、Mestel 円盤、回転曲線 & 第1章の数値積分の応用、格子収束 \\
付録B\newline Notebook B & §6.5\newline & 軸対称密度・ポテンシャル、Miyamoto--Nagai モデル & 2次元有限差分、対称軸の処理、疎行列、格子収束 \\
付録C\newline Notebook C & §7.2--7.3 & 軸対称・非軸対称ポテンシャル、エピサイクル運動、box・loop 軌道 & 第3章の時間積分法の応用、保存量による検証 \\
付録D & 直接の対応なし & 第2章・付録Bで現れる連立一次方程式の背景 & Gauss の消去法、LU 分解、疎行列、Jacobi 法 \\
\bottomrule
\end{tabularx}
\end{center}
(注) 表の「Rydenテキストとの対応」は、主として物理的な題材の対応を示しており、Rydenテキストが有限差分法や Poisson 方程式の数値ソルバーを扱っているという意味ではない。
}

\section*{本書で扱う範囲}

本書で学ぶ数値計算法は、次の4つである。

\begin{itemize}
  \item \textbf{数値積分}（第1章、付録A）
  \item \textbf{数値微分}（第2章）
  \item \textbf{有限差分法による境界値問題の解法}（第2章、付録B）
  \item \textbf{初期値問題に対する時間積分法}（第3章、付録C）
\end{itemize}

一方、補間法そのもの（スプライン補間など）、非線形方程式の解法（Newton 法など）、固有値問題は扱わない。
連立一次方程式はライブラリのソルバーで解き、その基本的な考え方を付録Dで紹介するにとどめる。
また、流体の方程式のような偏微分方程式を実際に時間発展させる数値解法そのものは扱わない。
ただし、境界条件、陽的・陰的な時間積分、粒子の初期位置分布や重力計算、および宇宙流体計算で使われる2次・3次精度 Runge--Kutta 法（RK2・RK3）との関係は、脚注や補足で紹介する。
これらは今後の学習への案内であり、本書で新たに実装する数値実験ではない。
これらの内容をさらに学ぶ場合は、巻末の参考図書を参照してほしい。

\section*{次の学習に向けて}

本教材で扱った数値積分、数値微分、有限差分法、時間積分法は、宇宙流体シミュレーションや重力多体系シミュレーションで使われる数値計算法の基礎となる。
一方、実際のシミュレーションでは、保存則に基づく離散化、数値流束、重力場の効率的な計算、複数の物理過程の結合など、さらに多くの数値計算法を組み合わせる必要がある。
本教材のあとに学ぶ内容は、目的に応じて次のように選ぶとよい。

\begin{itemize}
  \item \textbf{宇宙流体シミュレーションへ進む場合：}
  福江純・和田桂一・梅村雅之『宇宙流体力学の基礎［改訂版］』\cite{fukue2022}では、宇宙流体の基礎方程式、自己重力流体、波動、衝撃波、流体不安定性などを通して、宇宙流体力学そのものの物理的な基礎を学ぶことができる。
  富阪幸治・花輪知幸・牧野淳一郎編『シミュレーション天文学』\cite{tomisaka2007}では、常微分方程式と偏微分方程式の数値解法から、流体力学・磁気流体力学、自己重力を含む流体計算までを、天文シミュレーションとの関係を意識しながら学ぶことができる。
  藤井孝藏・立川智章『Pythonで学ぶ流体力学の数値計算法』\cite{fujii2020}では、圧縮性流体の数値計算法を Python による実装とともに学ぶことができる。
  松本倫明・大須賀健・須佐元『輻射電磁流体シミュレーションの基礎』\cite{matsumoto}では、有限体積法、数値流束、SPH 法、磁気流体力学、輻射流体力学など、宇宙物理シミュレーションで用いられる数値計算法をより詳しく学ぶことができる。

  \item \textbf{重力多体系・銀河力学シミュレーションへ進む場合：}
  James Binney \& Scott Tremaine, ``\textit{Galactic Dynamics}''\cite{bt2008}では、銀河の重力場の表現、多重極展開、基底関数展開、$N$ 体計算における重力計算法などを詳しく学ぶことができる。
  富阪幸治・花輪知幸・牧野淳一郎編『シミュレーション天文学』\cite{tomisaka2007}では、無衝突系・衝突系の重力多体系計算を含め、天文シミュレーションにおける $N$ 体計算を広く学ぶことができる。
  より専門的には、
  Sverre J. Aarseth, ``\textit{Gravitational N-Body Simulations: Tools and Algorithms}''\cite{aarseth2003}では、直接重力計算、時間積分、近接遭遇、正則化、個別時間刻み、精度と計算効率など、重力 $N$ 体シミュレーションのアルゴリズムを体系的に学ぶことができる。特に、星団のような衝突系の計算を詳しく学ぶ際の参考になる。
  Walter Dehnen \& Justin I. Read, ``\textit{N-body simulations of gravitational dynamics}''\cite{dehnen2011}は、衝突系と無衝突系の両方について、重力 $N$ 体シミュレーションで用いられる代表的な数値手法を整理したレビューであり、重力多体系計算の全体像を把握するための参考になる。
\end{itemize}

これらの専門的な教科書やレビューを読むときにも、本教材で繰り返した解析解との比較、収束試験、保存量による診断、誤差の原因を区別する考え方は、そのまま役立つ。
